\clearpage\pagestyle{otMain}\setlength{\otBlocEspacement}{2.5mm}\fontsize{11}{15}\selectfont\corrigesChapitre{21}{Les probabilités : a chances sur b}
\begin{blocLivre}{cor-21-1}\begin{corrige}{1}{Premières fractions de chances}\demarche{Compter les issues favorables sur le total.}
Obtenir 6 : $\res{1\ \text{chance sur }6}$. Pair : 3 issues, $\res{3\ \text{chances sur }6}$. De 1 à 6 : toutes les issues, $\res{6\ \text{chances sur }6}$ : certain.
\end{corrige}\end{blocLivre}
\begin{blocLivre}{cor-21-2}\begin{corrige}{2}{Le sac de boules}\demarche{Prendre les 10 boules comme issues.}
Rouge : $\res{2\ \text{chances sur }10}$ ; verte : $\res{3\ \text{sur }10}$ ; noire : $\res{5\ \text{sur }10}$, une chance sur deux.
\label{fig-cor-21-2}\figureGeo[corrige]{pro-2}
\end{corrige}\end{blocLivre}
\begin{blocLivre}{cor-21-3}\begin{corrige}{3}{Le jeu de 32 cartes}\demarche{Compter les cartes favorables.}
Cœur : $\res{8\ \text{chances sur }32}$. Rouge : $\res{16\ \text{sur }32}$, une chance sur deux. L’as de pique : $\res{1\ \text{sur }32}$.
\end{corrige}\end{blocLivre}
\begin{blocLivre}{cor-21-4}\begin{corrige}{4}{Équiprobable ou pas ?}\demarche{Vérifier l’égalité des secteurs avant de compter.}
Sur la roue $\res A$ : ses 8 secteurs identiques ont les mêmes chances. Sur B, le grand secteur sort plus souvent : les issues ne sont pas équiprobables, le compte « 1 sur 4 » serait faux.
\label{fig-cor-21-4}\figureGeo[corrige]{pro-4}
\end{corrige}\end{blocLivre}
\begin{blocLivre}{cor-21-5}\begin{corrige}{5}{Comparer deux évènements}\demarche{Lister les issues favorables des deux côtés.}
Plus grand que 4 : issues 5 et 6, soit 2 chances sur 6. Au plus 2 : issues 1 et 2, encore $\res{2\ \text{chances sur }6}$ : $\res{\text{égalité}}$ parfaite, contre l’intuition.
\end{corrige}\end{blocLivre}
\begin{blocLivre}{cor-21-6}\begin{corrige}{6}{La mémoire du dé}\demarche{Revenir à l’indépendance des lancers.}
Le dé $\res{\text{ne se souvient pas}}$ : chaque lancer est une nouvelle expérience, le 6 garde $\res{1\ \text{chance sur }6}$. Les lancers précédents n’influencent pas le suivant.
\end{corrige}\end{blocLivre}
\begin{blocLivre}{cor-21-7}\begin{corrige}{7}{Les deux pièces}\demarche{Distinguer PF et FP.}
Issues : PP, PF, FP, FF - quatre chemins équiprobables. Exactement un pile : PF et FP, soit $\res{2\ \text{chances sur }4}$, une chance sur deux.
\label{fig-cor-21-7}\figureGeo[corrige]{pro-7}
\end{corrige}\end{blocLivre}
\begin{blocLivre}{cor-21-8}\begin{corrige}{8}{Le tableau des deux dés}\demarche{Compter les cases du tableau.}
Le tableau $6\times6$ recense $\res{36}$ issues. La somme 7 sort par (1 ; 6), (2 ; 5), (3 ; 4), (4 ; 3), (5 ; 2), (6 ; 1) : $\res{6\ \text{chances sur }36}$.
\label{fig-cor-21-8}\figureGeo[corrige]{pro-8}
\end{corrige}\end{blocLivre}
\begin{blocLivre}{cor-21-9}\begin{corrige}{9}{Pile puis dé}\demarche{Multiplier les possibilités des deux étapes.}
$2\times6=\res{12}$ issues équiprobables. « Face puis 6 » est un seul chemin : $\res{1\ \text{chance sur }12}$.
\end{corrige}\end{blocLivre}
\begin{blocLivre}{cor-21-10}\begin{corrige}{10}{Plus probable que son contraire ?}\demarche{Comparer à la moitié du total.}
Gagnant : $\res{7\ \text{chances sur }10}$ ; perdant : $\res{3\ \text{sur }10}$. Le gagnant dépasse la chance sur deux : $\res{\text{oui}}$, il est plus probable que son contraire.
\end{corrige}\end{blocLivre}
\begin{blocLivre}{cor-21-11}\begin{corrige}{11}{Les deux tirages de jetons}\demarche{Remettre le jeton rend les étapes identiques.}
Issues : RR, RB, BR, BB - quatre chemins. Couleurs différentes : RB et BR, $\res{2\ \text{chances sur }4}$, une chance sur deux.
\label{fig-cor-21-11}\figureGeo[corrige]{pro-11}
\end{corrige}\end{blocLivre}
\begin{blocLivre}{cor-21-12}\begin{corrige}{12}{La loterie honnête ?}\demarche{Écrire « une chance sur 20 » sur le même total.}
$8$ sur $200$ se compare à $10$ sur $200$ (une chance sur 20) : $8<10$, donc $\res{\text{oui}}$, moins d’une chance sur 20. Les fractions de chances se comparent comme les fractions.
\end{corrige}\end{blocLivre}
\begin{blocLivre}{cor-21-13}\begin{corrige}{13}{Construire l’expérience}\demarche{Choisir 8 issues équiprobables dont 3 favorables.}
Par exemple : une roue à $\res8$ secteurs identiques dont $\res3$ marqués « gagné » ; ou un sac de 8 boules indiscernables dont 3 rouges. L’essentiel : 8 issues de mêmes chances, 3 favorables.
\end{corrige}\end{blocLivre}
\begin{blocLivre}{cor-21-14}\begin{corrige}{14}{Le détour par le contraire}\demarche{Compter ce qui reste quand l’évènement échoue.}
Ne pas obtenir 6 : les issues 1 à 5, $\res{5\ \text{chances sur }6}$. Avec « obtenir 6 » (1 sur 6), les deux comptes réunissent les $\res6$ issues : un évènement et son contraire se partagent toutes les issues.
\end{corrige}\end{blocLivre}
\begin{blocLivre}{cor-21-15}\begin{corrige}{15}{Le grand tri}\demarche{Compter les cartes favorables de chacun.}
1 sur 32, puis 4 sur 32, puis 8 sur 32, puis 32 sur 32 : $\res{\text{as de cœur}<\text{as}<\text{cœur}<\text{certain}}$. Le dernier évènement est réalisé par toutes les issues.
\end{corrige}\end{blocLivre}
